\section*{Chapitre \ref{ch:dv:jumps-and-levy-models} --- Sauts et modèles de Lévy}

\begin{solution}{exo:dv:jumps-and-levy-models:1}
Le nombre de sauts suit une loi de Poisson de moyenne $\lambda T$~: $1-e^{-3/52}=5.6\%$ sur une semaine, $1-e^{-3}=95.0\%$ sur un an.
\end{solution}

\begin{solution}{exo:dv:jumps-and-levy-models:2}
Pour $u=-i$, $iu=1$ et $u^2=-1$, donc l'exposant vaut
$T\bigl(-\frac12\sigma^2(1-1)-\lambda\bar k+\lambda(e^{\mu_J+\frac12\delta^2}-1)\bigr)=0$. $\phi(-i)=\E[e^{x_T}]=\E[S_T]/F_T$~:
le prix à terme du modèle est égal à celui du marché~: c'est la condition de martingale
que la correction de dérive $-\lambda\bar k$ impose.
\end{solution}

\begin{solution}{exo:dv:jumps-and-levy-models:3}
L'asymétrie varie comme $T^{-1/2}$~: $-0.4\sqrt{52}=-2.88$ à une semaine et $-0.4/\sqrt4=-0.20$ à quatre ans.
\end{solution}

\begin{solution}{exo:dv:jumps-and-levy-models:4}
Détenir $\Delta$ actions contre le straddle vendu~; le P\&L dans chaque état vaut $-(V_i-V_0)+\Delta(S_i-S_0)$. En égalant
les deux (une position sans risque à travers l'événement), on obtient $\Delta^*=(V_u-V_d)/(S_u-S_d)$. Ici
$(2.21+1.48)/(50\times(0.221+0.148))=3.69/18.45=0.20$.
\end{solution}

\begin{solution}{exo:dv:jumps-and-levy-models:5}
Le delta du modèle est la pente de la valeur du straddle pour un petit mouvement du spot avant l'événement, là où le
straddle est presque symétrique~: $-0.02$. L'événement n'est pas un petit mouvement~; le straddle gagne dans les deux issues, mais
davantage sur le mouvement haussier, donc la bonne couverture est la pente de la corde $+0.20$. Une position courte de 0.02 action
s'ajoute à l'exposition sur le mouvement haussier au lieu de la réduire.
\end{solution}

\begin{solution}{exo:dv:jumps-and-levy-models:6}
Dans un modèle de Lévy, l'asymétrie décroît comme $T^{-1/2}$ et le smile de volatilité s'aplatit rapidement, si bien qu'une calibration à un mois laisse
le skew à un an proche de zéro ($-0.07$ contre $-0.25$). Le skew de long terme du marché provient de la persistance et
de la corrélation de la volatilité, que Bates ajoute par le facteur de variance de Heston~; ses sauts maintiennent pentu le court terme.
\end{solution}

\begin{solution}{exo:dv:jumps-and-levy-models:7}
L'écart-type passe de 0.97 à 0.55 primes, et la perte centennale de 4.1 à 2.1~: des sauts plus petits
et plus fréquents se rapprochent davantage d'une diffusion. (La variance totale baisse aussi, à 13.5~\% par an, et la prime
à 1.55.)
\end{solution}

\begin{solution}{exo:dv:jumps-and-levy-models:8}
La majeure partie des 160~\% vient d'un seul saut le jeudi, et non d'une diffusion. À travers le saut, le delta de Black laisse un P\&L
d'écart-type d'environ 11~\% de la prime même avec des tailles d'issue connues, et une couverture ajustée au saut laisse 16~\%
une fois les tailles incertaines. Le risque porte sur la taille et le sens du mouvement, que seules les options (ou le fait de ne pas
être vendeur de l'événement) peuvent couvrir.
\end{solution}

\begin{solution}{pb:dv:jumps-and-levy-models:1}
\textbf{1.} 159.5~\% à la monnaie, 79.3~\% pour le strike 38 et 116.4~\% pour le strike 62.
\textbf{2.} La \omterm{def:dv:implied-volatility-and-its-surface:density}{densité risque-neutre} est bimodale~: deux bosses autour des deux issues, avec peu de masse exactement à la monnaie. Les options
proches de la monnaie sont chères par rapport à une lognormale aux mêmes ailes, ce qui dessine un froncement.
\textbf{3.} $1.595^2/52-0.35^2/52=0.0465$~; un écart-type de 21.6~\%.
\textbf{4.} $21.6\%\times\sqrt{2/\pi}=17.2\%$. Le straddle paie la variation absolue, si bien que son prix mesure
l'espérance de ce mouvement absolu, quelle que soit la distribution~; l'écart-type s'interprète à travers une hypothèse normale que
un mouvement binaire contredit.
\textbf{5.} La distribution n'est pas lognormale~: une volatilité ajuste le prix d'un seul strike, et non la forme.
\textbf{6.} $\sigma=40.6\%$, $p=0.399$, $a=0.199$, avec une erreur de 0.07 point de volatilité.
\textbf{7.} $+22.1\%$ et $-14.7\%$.
\textbf{8.} L'incertitude sur la taille des issues élargit chaque bosse~; le modèle à deux points ne peut les élargir que
par la diffusion.
\textbf{9.} Environ 40~\%~: le prix de marché du succès. Un trader le compare à sa propre estimation (issue de données d'essais,
précédents) et négocie la différence avec des options digitales ou des call spreads.
\textbf{10.} $(1-0.4\times1.2214)/0.6-1=-14.8\%$.
\textbf{11.} 8.86~; 1.81.
\textbf{12.} $+0.09$ et $-0.02$~; 1.01 et 2.01.
\textbf{13.} $0.20$~; zéro.
\textbf{14.} $0.20$~; 1.45, 16.5~\% de la prime.
\textbf{15.} Options sur l'événement~: racheter une partie du straddle, ou des strangles et des papillons qui portent
l'exposition à la taille des issues.
\textbf{16.} Le delta est la réponse aux petits mouvements~; le ratio est la pente de la corde entre les deux issues. Ils ne coïncident que
pour un profil de gains linéaire.
\textbf{17.} Couvrir la diffusion en delta normalement jusqu'à l'événement~; juste avant, passer au ratio de corde pour les
deux issues~; après l'événement, revenir au delta.
\textbf{18.} Au moins le résidu de la meilleure couverture, environ 1.45 par straddle (un écart-type), davantage pour un
niveau de confiance~; à multiplier par la position et agrégé sur les événements.
\textbf{19.} Une probabilité de succès de 0.40 avec des mouvements de $+22.1\%$ et $-14.7\%$~; le delta de Black laisse un P\&L
d'écart-type de 1.01 (11.4~\% de la prime) à travers l'événement, le ratio binomial 0.20 l'annule quand les tailles
sont connues, et laisse 1.45 (16.5~\%) quand elles sont incertaines.
\textbf{20.} La distribution du mouvement, et avec elle le ratio de couverture à travers le saut,
qu'aucune volatilité unique ne contient.
\end{solution}

\begin{solution}{iq:dv:jumps-and-levy-models:1}
Dans une diffusion, une option hors de la monnaie donnée exige un mouvement de nombreux écarts-types quand $T\to0$, si bien que son prix
tend vers zéro exponentiellement vite~; avec des sauts, un seul saut suffit et le prix est proportionnel à $T$. Les prix de marché des
options hors de la monnaie à courte échéance diminuent à peu près linéairement, et les smiles à courte échéance sont pentus.
\iqlookfor{les vitesses de convergence, et l'idée de Carr--Wu.}
\end{solution}

\begin{solution}{iq:dv:jumps-and-levy-models:2}
$dS/S=-\lambda\bar k\,dt+\sigma\,dW+(e^J-1)\,dN$, $J\sim\mathcal N(\mu_J,\delta^2)$, $N$ de loi Poisson($\lambda$). Étant donné $n$
sauts, le log-rendement est normal, si bien que $C=\sum_n e^{-\lambda T}(\lambda T)^n/n!\;C_{\mathrm{BS}}(F_n,K,T,\sigma_n)$ avec
$F_n=Fe^{n(\mu_J+\delta^2/2)-\lambda\bar kT}$ et $\sigma_n^2=\sigma^2+n\delta^2/T$.
\iqlookfor{le compensateur, et l'argument de conditionnement.}
\end{solution}

\begin{solution}{iq:dv:jumps-and-levy-models:3}
Non~: à travers un saut, la position couverte rapporte $V(Se^J)-V(S)-\Delta S(e^J-1)$, ce qu'aucun $\Delta$ unique n'annule pour
tout $J$. Avec deux issues connues, un ratio de corde le fait~; sinon, couvrir la convexité résiduelle avec des options (elles ont
la même exposition au saut), dimensionner les positions sur le résiduel, et conserver une réserve.
\iqlookfor{la convexité à travers le saut, et les options comme couverture.}
\end{solution}

\begin{solution}{iq:dv:jumps-and-levy-models:4}
Les cumulants sont linéaires en $T$, si bien que l'asymétrie $\propto T^{-1/2}$ et le kurtosis en excès $\propto T^{-1}$~: des smiles pentus à
courte échéance, qui s'aplatissent vite. Les smiles des indices s'aplatissent plus lentement, ce qui appelle aussi un \omterm{def:dv:stochastic-volatility:sv}{modèle à volatilité stochastique}.
\iqlookfor{la loi d'échelle, et sa conséquence.}
\end{solution}

\begin{solution}{iq:dv:jumps-and-levy-models:5}
La fonction caractéristique du \omterm{def:dv:jumps-and-levy-models:vg}{modèle variance gamma} ne décroît que comme $|u|^{-2T/\nu}$, très lentement à une semaine, alors que celle de Heston
décroît exponentiellement~; une grille tronquée à un $u$ modéré perd une grande partie de la queue. Élargir la grille (ou utiliser un schéma amorti ou
corrigé en queue), et tester par rapport à la forme fermée conditionnelle à l'échéance la plus courte.
\iqlookfor{le taux de décroissance de la fonction caractéristique.}
\end{solution}

\begin{solution}{iq:dv:jumps-and-levy-models:6}
Séparer la diffusion de l'événement~: évaluer avec une diffusion à la volatilité ordinaire plus un saut dont la loi se lit
à partir du smile (une \omterm{def:dv:parametrising-the-surface:event}{variance d'événement}, ou une calibration à deux points si le smile présente un froncement). Couvrir la diffusion en delta
avant l'événement et le saut avec le ratio de corde ou avec des options~; après l'événement, la volatilité implicite se replie
sur celle de la diffusion. Ce qui peut mal tourner~: la taille ou le sens du mouvement, un mouvement plus grand que celui intégré dans les prix, la liquidité à l'ouverture,
et l'effondrement de la volatilité des autres échéances.
\iqlookfor{diffusion plus événement, la couverture à travers le saut, l'effondrement de la volatilité.}
\end{solution}
